

\Needspace{90mm}
\TypeQuestion{}{parallel_segment_ratio_multiple_cuts}
\FigProblemBlock{48mm}{fig/lesson_08_multiple_cuts}{
\Text{$\Tri{ABC}$ の辺 $\Lseg{AB}$ 上に点 $\mathrm{D},\mathrm{F}$、辺 $\Lseg{AC}$ 上に点 $\mathrm{E},\mathrm{G}$ を、それぞれ $\mathrm{A},\mathrm{D},\mathrm{F},\mathrm{B}$ および $\mathrm{A},\mathrm{E},\mathrm{G},\mathrm{C}$ の順にとる。$\Lseg{DE}\Parallel\Lseg{FG}\Parallel\Lseg{BC}$ であり、$\Lseg{AD}=4\mathrm{cm},\ \Lseg{DF}=2\mathrm{cm},\ \Lseg{FB}=3\mathrm{cm},\ \Lseg{AE}=6\mathrm{cm},\ \Lseg{BC}=18\mathrm{cm}$ とする。}
}{
\QQ{辺 $\Lseg{AC}$ の長さを求めなさい。}{}
\QQ{線分 $\Lseg{EG}$ の長さを求めなさい。}{}
\QQ{線分 $\Lseg{FG}$ の長さを求めなさい。}{}
}


\Needspace{90mm}
\TypeQuestion{}{parallel_segment_ratio_trapezoid_diagonals}
\FigProblemBlock{48mm}{fig/lesson_08_trapezoid}{
\Text{$\Lseg{AB}\Parallel\Lseg{CD}$ である台形 $\mathrm{ABCD}$ の対角線 $\Lseg{AC},\Lseg{BD}$ の交点を $\mathrm{O}$ とする。$\Lseg{AB}=6\mathrm{cm},\ \Lseg{CD}=10\mathrm{cm},\ \Lseg{AO}=9\mathrm{cm},\ \Lseg{BO}=6\mathrm{cm}$ とする。}
}{
\QQ{$\Lseg{AO}:\Lseg{OC}$ を求めなさい。}{}
\QQ{線分 $\Lseg{OD}$ の長さを求めなさい。}{}
\QQ{対角線 $\Lseg{AC}$ の長さを求めなさい。}{}
}

%補助線は自分で引かせる形にしてください。
\TypeQuestion{}{parallel_segment_ratio_auxiliary_line}
\FigProblemBlock{48mm}{fig/lesson_08_auxiliary}{
\Text{$\Lseg{AB}\Parallel\Lseg{CD}$ である台形 $\mathrm{ABCD}$ で、$\Lseg{AB}=7\mathrm{cm},\ \Lseg{CD}=19\mathrm{cm},\ \Lseg{AD}=15\mathrm{cm}$ とする。辺 $\Lseg{AD}$ 上に $\Lseg{AE}=5\mathrm{cm}$ となる点 $\mathrm{E}$ をとり、点 $\mathrm{E}$ を通り辺 $\Lseg{AB}$ に平行な直線と辺 $\Lseg{BC}$ との交点を $\mathrm{F}$ とする。次の補助線を引いて考えなさい。}
}{
\Text{点 $\mathrm{A}$ を通り辺 $\Lseg{BC}$ に平行な直線を引き、線分 $\Lseg{EF},\Lseg{CD}$ との交点をそれぞれ $\mathrm{G},\mathrm{H}$ とする。}
\QQ{線分 $\Lseg{DH}$ の長さを求めなさい。}{}
\QQ{線分 $\Lseg{EG}$ の長さを求めなさい。}{}
\QQ{線分 $\Lseg{EF}$ の長さを求めなさい。}{}
